\section{Discussion}
\label{sec:limits}
In the exploratory R13 application run, both model candidates fail the
configured eligibility and performance gates. Signal pairs were selected
after a pilot, and the exact runtime source was not pinned; the retained
correlations and NRMSE values therefore document the recorded gate path rather
than confirmatory model performance. In R20, six sequential tasks complete a
synthetically approved, readback-verified recipe trial and task closure. Each
has one run-tagged sample, collected after a 20-s wait initiated after task
completion and evidence queries through an isolated SQLite fallback; all six
samples remain outside their target bands. This supports the tested
application path, not successful process control.

R20 uses one simulator instance per scene and does not establish independent
replication, convergence, or quality improvement. Its post-task sample is not
a response trajectory or a fixed command-to-sample latency measurement.
Controlled profile/binding adaptation, cross-harness portability, normal
online DAQ backend behavior, and physical transfer remain untested. The
synthetic approval records do not measure operator effectiveness
\cite{parasuramanmanzey2010complacency,parasuraman2000model,bainbridge1983ironies};
restricted tool access does not establish deployment security
\cite{greshake2023injection,liu2024pinjection,iec62443-3-3} or replace
independently engineered PLC/SIS protection \cite{iec61511}.

Further evaluation needs a versioned build, matched configuration changes,
independent simulator replications, normal online DAQ, and physical-plant
trials.

\section{Conclusion}
\label{sec:conclusion}
AgentWorkShop provides an implemented task-to-device integration architecture
that joins collaborative harnesses, production context, operation checks,
semantic I/O, and route-specific records. R13 and R20 demonstrate selected
simulator paths: R13's model candidates fail their gates, and all six R20
samples remain outside target. The evidence establishes integration behavior
and does not establish process-control efficacy.
